\part{Multivariable Calculus}
\chapter{Differentiation in $\bR^n$}
\section{Differentials and Partial Derivatives}
\begin{importantdefinition}{Fréchet Derivative}{}
	Let $V,W$ be normed vector spaces. Let $U$ be an open subset of $V$. Then $f : U \to W$ is called \tib{ Fréchet differentiable} or just \tib{differentiable} at $p \in U$ if there exists a continuous linear map $Df(p) : V \to W$ such that
	\begin{align*}
		\lim_{h \to 0} \frac{\norm{f(p+h) - f(p) - Df(p)(h)}_W}{\norm{h}_V} = 0.
	\end{align*}
	We call $Df(p)$ the \tib{differential of $f$ at $p$}. We call $f$ \tib{differentiable} if it is differentiable at every point $p \in U$. We call $f$ \tib{continuously differentiable} if the map $Df : p \mapsto Df(p)$ is continuous.
\end{importantdefinition}
Recall that linear maps are continuous if and only if they are bounded, which is always the case if $V$ is finite-dimensional.
\begin{exercise}
	Verify that if we identify scalars $c \in \bR$ with linear maps $x \mapsto c \cdot x$, this definition agrees with the usual definition for $f : \bR \to \bR$.
\end{exercise}
\begin{theorem}
	If $f : U \to W$ is differentiable, it is continuous.
\end{theorem}
\begin{proof}
	If $f$ is differentiable, we have
	\begin{align*}
		\lim_{\norm{h}_V \to 0} \frac{\norm{f(p+h) - f(p) - Df(p)(h)}_W}{\norm{h}_V} = 0.
	\end{align*}
	This can only be the case if
	\begin{align*}
		\lim_{\norm{h}_V \to 0} \norm{f(p+h) - f(p) - Df(p)(h)}_W = 0,
	\end{align*}
	i.e. if
	\begin{align*}
		\lim_{\norm{h}_V \to 0} f(p+h) = \lim_{\norm{h}_V \to 0} f(p) + Df(p)(h).
	\end{align*}
	Since $Df(p)$ is continous with $Df(0) = 0$, this simplifies to
	\begin{align*}
		\lim_{\norm{h}_V \to 0} f(p+h) = f(p),
	\end{align*}
	which is the definition of sequential continuity.
\end{proof}
\begin{importantdefinition}{Directional and Partial Derivates}{}
	Let $V$ and $W$ be normed vector spaces. Let $v \in V$. Let Let $U \subseteq V$ be open and $f : U \to W$. Then we call the limit
	\begin{align*}
		\pderivative{v} f(p) = \frac{d}{dh} f(p + hv) \bigg |_{h = 0} = \lim_{h \to 0} \frac{f(p + hv) - f(p)}{h}
	\end{align*}
	the \tib{directional derivative} of $f$ at $p$ in direction $v$. If $(x_0, \hdots, x_n)$ is a basis of $V$, we call the directional derivatives
	\begin{align*}
		\partial_i f(p) := \pderivative{x_i} f(p) 
	\end{align*}
	the \tib{partial derivatives of $f$ at $p$}. If all partial derivatives exist, we call $f$ \tib{partially differentiable}. If all maps $\partial_i f : p \mapsto \partial_if(p)$ are continuous, we call $f$ \tib{continuously partially differentiable}.
\end{importantdefinition}
Partial derivatives aren't actually any more complicated than standard one-dimensional derivatives - you can calculate partial derivatives by simply fixing all coordinates except one and then taking a regular derivative in the remaining coordinate. For example, if we have a function
\begin{align*}
	f : \bR^3 &\to \bR\\
		(x,y,z) &\mapsto f(x,y,z),
\end{align*}
we can define a function $f|_1 : \bR \to \bR$ by restriction to the first coordinate:
\begin{align*}
	f|_1(x) := f(x,y,z).
\end{align*}
Using this function, we can simplify the partial derivative:
\begin{align*}
	\partial_1 f(x,y,z) 
	&= \lim_{h \to 0} \frac{f((x,y,z) + he_1) - f(x,y,z)}{h}\\
	&= \lim_{h \to 0} \frac{f(x+h,y,z) - f(x,y,z)}{h}\\
	&= \lim_{h \to 0} \frac{f|_1(x+h) - f(x)}{h}\\
	&= \frac{d}{dx} f|_1(x)
\end{align*}
This tells us that, for example, the function $f(x,y,z) = x \cdot y \cdot z$ has partial derivatives
\begin{align*}
	\partial_1 f(x,y,z) = \frac{d}{dx} xyz = yz\\
	\partial_2 f(x,y,z) = \frac{d}{dy} xyz = xz\\
	\partial_3 f(x,y,z) = \frac{d}{dz} xyz = xy.\\
\end{align*}
\begin{exercise}
	\begin{enumerate}
		\item Show that if $f$ is differentiable, then all directional derivatives exist and fulfill $\pderivative{v} f(p) = Df(p)(v)$.
		\item Conclude that $Df(p)$ is represented in basis $(x_0, \hdots, x_n)$ by the \tib{Jacobian matrix}
		\begin{align*}
			Jf(p)
			=
			\lr(\pderivative{x_1} f(p), \hdots, \pderivative{x_n} f(p))
			=
			\begin{pmatrix}
				\pderivative{x_1}f_1(p) & \hdots & \pderivative{x_n}f_1(p)\\
				\vdots && \vdots\\
				\pderivative{x_1}f_n(p) & \hdots & \pderivative{x_n}f_n(p)
			\end{pmatrix}
		\end{align*}
	
	\end{enumerate}
\end{exercise}
If the image of $f$ is one-dimensional, the Jacobian matrix is simply a row vector. We call the transpose of this vector the \tib{gradient of $f$ at $p$} and denote it
\begin{align*}
	\nabla f(p)
	:=
	Jf(p)^\top
	=
	\begin{pmatrix}
		\pderivative{x_1} f(p)\\
		\vdots
		\pderivative{x_n} f(p)
	\end{pmatrix}
\end{align*}
The transposition comes from the fact that we would like to be able to get the directional derivatives of $f$ at $p$ by taking a scalar product with the gradient, which gives us
\begin{align*}
	\scalar{\nabla f(p)}{v} 
	&= \nabla f(p)^\top v\\ 
	&= Jf(p) \cdot v\\ 
	&= Df(p)(v)\\
	&= \pderivative{v} f(p)
\end{align*}
There is one very important pitfall to watch out for when it comes to differentiability: We have just shown that if $f$ is differentiable, we can get the differential by construction the Jacobian matrix from the partial derivatives. This makes it look like every partially differentiable function is differentiable, which turns out to not be true in general. 
\newpar
For example, there exist functions which are partially differentiable, but not continuous - we showed earlier that all differentiable functions are continuous, so these functions are quick counterexamples.
\newpar
One example is
\begin{align*}
	f : \bR^2 \to \bR\\
	(x,y) &\mapsto
	\begin{cases}
		\frac{xy}{(x^2 + y^2)^2} & (x,y) \neq 0\\
		0 & (x,y) = 0
	\end{cases}.
\end{align*}
This function has $\partial_1 f(0) = \partial_2 f(0) = 0$, since $f(h,0) = f(0,h) = 0$ for all $h \in \bR$, but isn't continuous, since $\lim_{n \to \infty} f\lr(\frac{1}{n}, \frac{1}{n}) = \lim_{n \to \infty} \frac{n^2}{4} = \infty \neq f(0)$.
\newpar
For a more involved counterexample, consider the function
\begin{align*}
	f : \bR^2 &\to \bR\\
	(x,y) &\mapsto 
	\begin{cases}
		\frac{y^3}{x^2 + y^2} & (x,y) \neq 0\\
		0 & (x,y) = 0
	\end{cases}
\end{align*}
This function is continuous - for $(x,y) \neq 0$, it is a composition of continuous functions, and thus continuous, and since $x^2 + y^2 \geq y^2$, we additionally have
\begin{align*}
	\abs{f(x,y)} 
	&= \frac{\abs{y^3}}{\abs{x^2 + y^2}}\\
	&= \frac{\abs{y^3}}{x^2 + y^2}\\
	&\leq \frac{\abs{y^3}}{y^2}\\
	&= \abs{y},
\end{align*}
so $f(x,y)$ converges to $f(0,0) = 0$ as $(x,y)$ converges to $0$, implying that $f$ is also continuous at the origin.
\newpar
$f$ is also directionally differentiable in all directions:
\begin{align*}
	\partial_v f(0)
	&= \lim_{h \to 0} \frac{f(h \cdot v) - f(0)}{h}\\
	&= \lim_{h \to 0} \frac{f(h \cdot v)}{h}\\
	&= \lim_{h \to 0} \frac{\lr(\frac{h^3v_2^3}{h^2 v_1^2 + h^2 v_2^2})}{h}\\
	&= \lim_{h \to 0} \frac{h^3v_2^3}{h^3 v_1^2 + h^3 v_2^2}\\
	&= \lim_{h \to 0} \frac{v_2^3}{v_1^2 + v_2^2}\\
	&= \frac{v_2^3}{v_1^2 + v_2^2}\\ 
	&= f(v)
\end{align*}
In particular, $\partial_1 f(0) = 0$ and $\partial_2 f(0) = 1$. However, $f$ is not differentiable at $0$, since otherwise our representation $Df(0) = Jf(0)$ would give us:
\begin{align*}
	\lim_{\vh \to 0} \frac{\abs{f(h) - f(0) - Df(0)(\vh)}}{\norm{\vh}_2}
	&=
	\lim_{\vh \to 0} \abs{\frac{f(h) - Df(0)(h)}{\norm{\vh}_2}}\\
	=
	\lim_{\vh \to 0} \abs{\frac{f(h) - \partial_1 f(0) \cdot h_1 + \partial_2 f(0) \cdot h_2}{\sqrt{h_1^2 + h_2^2}}}
	&=
	\lim_{\vh \to 0} \abs{\frac{f(h) - h_2}{\sqrt{h_1^2 + h_2^2}}}\\
	=
	\lim_{\vh \to 0} \abs{\frac{1}{\sqrt{h_1^2 + h_2^2}} \cdot \lr(\frac{h_2^3}{h_1^2 + h_2^2} - h_2)}
	&=
	\lim_{\vh \to 0} \abs{\frac{1}{\sqrt{h_1^2 + h_2^2}} \cdot \lr(\frac{h_2^3 - h_2(h_1^2 + h_2^2)}{h_1^2 + h_2^2})}\\
	=
	\lim_{\vh \to 0} \abs{\frac{1}{\sqrt{h_1^2 + h_2^2}} \cdot \lr(\frac{h_2^3 - h_2h_1^2 - h_2^3}{h_1^2 + h_2^2})}
	&=
	\lim_{\vh \to 0} \abs{\frac{1}{\sqrt{h_1^2 + h_2^2}} \cdot \lr(\frac{h_2h_1^2}{h_1^2 + h_2^2})}
	=
	\lim_{\vh \to 0} \abs{\frac{h_2h_1^2}{\sqrt{h_1^2 + h_2^2}^3}}
\end{align*}
Since we are only interested in whether this expression converges to $0$, we can square it to simplify slightly:
\begin{align*}
	\abs{\frac{h_2h_1^2}{\sqrt{h_1^2 + h_2^2}^3}}^2
	&=
	\abs{\frac{h_2^2h_1^4}{(h_1^2 + h_2^2)^3}}
\end{align*}
If we now set $h := \lr(\frac{1}{n}, \frac{1}{n})$, we get:
\begin{align*}
	\abs{\frac{h_2^2h_1^4}{(h_1^2 + h_2^2)^3}}
	&=
	\abs{\frac{\frac{1}{n^6}}{\lr(\frac{2}{n^2})^3}}
	=
	\abs{\frac{\frac{1}{n^6}}{\frac{8}{n^6}}}
	=
	\frac{1}{8}
	\neq 
	0
\end{align*}
The important takeaway here is that directional differentiability only guarantees "differentiability along straight paths", while full differentiability also requires "differentiability along curved paths".
\begin{importanttheorem}{}{}
	$f : V \to W$ is continuously differentiable if and only if it is continuously partially differentiable.
\end{importanttheorem}
\begin{proof}
	\begin{enumerate}
		\item[$\Rightarrow$:] We have $\Hom(\bR^n, \bR^m) \cong \bR^{mn}$. Therefore, the map $Jf : \bR^n \mapsto \bR^{mn}$ is continuous if and only if every component $(Jf)_i : \bR^n \mapsto \bR$ is continuous. By the definition of $Jf$, these are the partial derivatives.
		\item[$\Leftarrow$:]
		Since every normed complex vector space is isomorphic to a normed real vector space, we can assume without loss of generality that $W$ is real. Since all norms are equivalent, and since a sequence in $W$ converges iff. all components converge, we can additionally assume that $W$ is one-dimensional, i.e. $W \cong \bR$.
		\newpar
		We need to show that
		\begin{align*}
			\lim_{h \to \vzero} \frac{f(p+h) - f(p) - Jf(p)h}{\norm{h}_2} = \vzero
		\end{align*}
		Let $h := \sum_{i = 1}^n h_ie_i$. We decompose the path from $p$ to $p+h$ into its partial steps, taking one coordinate direction at a time:
		\begin{align*}
			\forall &k \in [0..n]:\\
			&p_k = p + \sum_{i = 1}^k h_i e_i
		\end{align*}
		Now, by the mean value theorem, there exists an $s_k \in [0,1]$ such that
		\begin{align*}
			f(p_k) - f\lr(p_{k-1})
			&=
			f(p_{k-1} + h_ke_k) - f\lr(p_{k-1})\\
			&= \partial_k f(p_{k-1} - s_kh_ne_n)h_k
		\end{align*}
		Now, since 
		\begin{align*}
			Jf(p)(h) = \sum_{i = 1}^n \partial_i f(p)h_i,
		\end{align*}
		we get:
		\begin{align*}
			&\frac{\abs{f(p+h) - f(p) - Jf(p)(h)}}{\norm{h}_2}\\
			=
			&\frac{1}{\norm{h}_2}\abs{f(p_n) - f(p_0) - \sum_{i = 1}^n \partial_i f(p)h_i}\\
			=
			&\frac{1}{\norm{h}_2}
			\abs{\sum_{i = 1}^n f(p_i) - f(p_{i-1}) - \partial_i f(p) h_i}\\
			=
			&\frac{1}{\norm{h}_2}
			\abs{\sum_{i = 1}^n \partial_i f(p_{i-1} + s_ih_ie_i)h_i - \partial_i f(p) h_i}\\
			=
			&\frac{1}{\norm{h}_2}
			\abs{\sum_{i = 1}^n \lr(\partial_i f(p_{i-1} + s_ih_ie_i) - \partial_i f(p)) h_i}\\
			=
			&\frac{1}{\norm{h}_2}
			\abs{\sum_{i = 1}^n \lr(\partial_i f\lr(p + \sum_{j = 1}^{i-1} h_je_j + s_ih_ie_i) - \partial_i f(p)) h_i}\\
			\leq
			&\frac{1}{\norm{h}_2}
			\sum_{i = 1}^n \abs{\lr(\partial_i f\lr(p + \sum_{j = 1}^{i-1} h_je_j + s_ih_ie_i) - \partial_i f(p)) h_i}\\
			=
			&
			\sum_{i = 1}^n \abs{\partial_i f\lr(p + \sum_{j = 1}^{i-1} h_je_j + s_ih_ie_i) - \partial_i f(p)} \cdot \frac{\abs{h_i}}{\norm{h}_2}\\
			\leq
			&
			\sum_{i = 1}^n \abs{\partial_i f\lr(p + \sum_{j = 1}^{i-1} h_je_j + s_ih_ie_i) - \partial_i f(p)}
		\end{align*}
		Since $s_i \in [0,1]$, we have
		\begin{align*}
			\lim_{h \to 0} p + \sum_{j = 1}^{i-1} h_je_j + s_ih_ie_i = p
		\end{align*}
		And thus, since the partial derivatives $\partial_i f$ are continuous, the right side converges to zero, so $f$ is differentiable.
	\end{enumerate}
\end{proof}
\begin{exercise}
	Let $V$ be an inner product space over a field $\bF = \bR$ or $\bF = \bC$. Let $v \in V$. Then the function 
	\begin{align*}
		\scalar{\cdot}{v} : V &\to \bF\\
		w &\mapsto \scalar{w}{v}
	\end{align*}
	is differentiable, with $D\scalar{\cdot}{v}(w) = v^\top$ for all $w \in V$.
\end{exercise}
\begin{exercise}
	Let $b$ be a symmetric bilinear form. Let
	\begin{align*}
		q(v) := b(v,v)
	\end{align*}
	be the corresponding quadratic form. Then
	$q$ is differentiable, with
	\begin{align*}
		Dq(v)(w) = 2b(v,w)
	\end{align*}
\end{exercise}
\clearpage

\section{Rules of Differentiation}
\begin{exercise}
	Show that the differential operator $D$ is linear, i.e. that we have
	\begin{align*}
		D(\alpha f + \beta g)(p) = \alpha Df(p) + \beta Dg(p)
	\end{align*}
	for $f, g$ differentiable at $p$.
\end{exercise}
\begin{importanttheorem}{Chain Rule in $\bR^n$}{}
	Let $f : U \to \bR^m$ and $g : V \to \bR^r$ with $U \subset \bR^n$, $V \subset \bR^m$ open and $f(U) \subseteq V$. Let $f$ be differentiable at a point $p \in U$ and let $g$ be differentiable at $f(p)$. Then
	\begin{align*}
		D(g \circ f)(p)
		&=
		Dg(f(p)) \circ Df(p)
	\end{align*}
\end{importanttheorem}
\begin{exercise}
	Verify that this version of the chain rule agrees with the familiar chain rule for $f,g : \bR \to \bR$.
\end{exercise}
\begin{proof}
	By definition of the differential, we have:
	\begin{align*}
		f(x) 
		&= f(p)\\
		&\+ 
		Df(p)(x-p)\\
		&\+ 
		\epsilon_f(x)\norm{x - p}\\
		g(f(x))
		&=
		g(f(p))\\
		&\+
		Dg(f(p))(f(x)-f(p))\\
		&\+
		\epsilon_g(f(x))\norm{f(x) - f(p)}
	\end{align*}
	For some functions $\epsilon_f : U \to \bR^m$, $\epsilon_g : U \to \bR^r$ which are continuous in $p$ and $f(p)$ and fulfill $\epsilon_f(p) = 0$ and $\epsilon_g(f(p)) = 0$. 
	
	The second equation tells us:
	\begin{align*}
		g(f(x))
		&=
		g(f(p))\\
		&\+
		Dg(f(p))(f(x) - f(p))\\
		&\+ \epsilon_g(f(x))\norm{f(x) - f(p)}
	\end{align*}
	Plugging in our equation for $f(x)$ now lets us regroup our terms into an equation of the same form as the definition of the differential:
	\begin{align*}
		g(f(x))
		&=
		g(f(p))\\
		&\+
		\tc{c3}{Dg(f(p))}
		\biggl[
			\overbrace{
			\tc{c2}{f(p)}
			+ 
			\tc{c4}{Df(p)(x-p)}
			+ 
			\epsilon_f(x)\tc{c1}{\norm{x - p}}}^{f(x)} \tc{c2}{- f(p)}
		\biggr]
		\\
		&\+ 
		\epsilon_g(f(x))
		\norm{
			\overbrace{\tc{c2}{f(p)}
			+ 
			\tc{c4}{Df(p)(x-p)}
			+ 
			\epsilon_f(x)\tc{c1}{\norm{x - p}}}^{f(x)} \tc{c2}{- f(p)}
		}
		\\
	\\
		&=
		g(f(p))\\
		&\+
		\tc{c3}{Dg(f(p))}
		\biggl[
			\tc{c4}{Df(p)(x-p)}
			+ 
			\epsilon_f(x)\tc{c1}{\norm{x - p}}
		\biggr]
		\\
		&\+ 
		\epsilon_g(f(x))
		\norm{
			\tc{c4}{Df(p)(x-p)}
			+ 
			\epsilon_f(x)\tc{c1}{\norm{x - p}}
		}
		\\
	\\
		&=
		g(f(p))\\
		&\+
		\tc{c3}{Dg(f(p))}\tc{c4}{Df(p)(x-p)}\\
		&\+ 
		\tc{c3}{Dg(f(p))}\epsilon_f(x)\tc{c1}{\norm{x - p}}
		\\
		&\+ 
		\epsilon_g(f(x))
		\norm{
			\tc{c4}{Df(p)(x-p)}
			+ 
			\epsilon_f(x)\tc{c1}{\norm{x - p}}
		}
		\\
	\\
		&=
		g(f(p))\\ 
		&\+ 
		\tc{c3}{Dg(f(p))}\tc{c4}{Df(p)(x-p)}\\
		&\+ 
		\underbrace{
			\lr(\epsilon_f(x)\tc{c3}{Dg(f(p))}
			+ 
			\epsilon_g(f(x))
			\norm{\tc{c4}{Df(p)} \lr(\frac{\tc{c4}{x-p}}{\tc{c1}{\norm{x-p}}}) + \epsilon_f(x)})
		}
		_
		{:= \epsilon_{g \circ f}(x)}
		\tc{c1}{\norm{x - p}}
	\\
		&=
		g(f(p))\\ 
		&\+ (\tc{c3}{Dg(f(p))} \circ \tc{c4}{Df(p)})\tc{c4}{(x-p)}\\
		&\+ \epsilon_{g \circ f}(x)\tc{c1}{\norm{x - p}} 
	\end{align*}
	All that remains is to show that $\epsilon_{g \circ f}(x)$ is continuous if we set
	\begin{align*}
		\epsilon_{g \circ f}(p) = 0.
	\end{align*}
	Since we have $\epsilon_f(p) = 0$ and $\epsilon_g(f(p)) = 0$, we have our desired result if we can show that 
	\begin{align*}
		\norm{Df(p) \lr(\frac{x - p}{\norm{x-p}})}
		&=
		\frac{\norm{Df(p)(x-p)}}{\norm{x-p}}.
	\end{align*}
	is bounded.
	Since $\bR^n$ is finite-dimensional, we can find a constant $c$ such that
	\begin{align*}
		\norm{Df(p)(x-p)} \leq c \cdot \norm{x-p},
	\end{align*}
	which means that our quotient is indeed bounded, completing the proof.
\end{proof}
\begin{importanttheorem}{Product Rule in $\bR^n$}{}
	Let $U \subseteq \bR^n$. Let $f,g : U \to \bR$ at a point $p \in U$. Then
	\begin{align*}
		D(f \cdot g)(p) = Df(p) \cdot g(p) + f(p) \cdot Dg(p).
	\end{align*}
\end{importanttheorem}
\begin{exercise}
	Verify that this version of the product rule agrees with the familiar product rule for $f,g : \bR \to \bR$.
\end{exercise}
\begin{proof}
	We define the multiplication function $m(x,y) := x \cdot y$. The partial derivatives are $\partial_1 m(x,y) = y$ and $\partial_2 m(x,y) = x$, i.e.
	\begin{align*}
		Dm(x,y) = 
		\begin{pmatrix}
			y & x\\
		\end{pmatrix}
	\end{align*}
	Applying the chain rule now gives
	\begin{align*}
		D(f \cdot g)(p)
		&=
		D(m \circ (f,g))(p)\\
		&=	
		Dm((f,g)(p)) \circ D(f,g)(p)\\
		&=	
		Dm(f(p), g(p)) \circ (Df(p), Dg(p))\\
		&=	
		Dm(f(p), g(p)) \cdot (Df(p), Dg(p))\\
		&=	
		\begin{pmatrix} 
			g(p) & f(p)
		\end{pmatrix}
		\cdot 
		\begin{pmatrix}
			Df(p)\\
			Dg(p)
		\end{pmatrix}\\
		&= Df(p) \cdot g(p) + f(p) \cdot Dg(p).
	\end{align*}
\end{proof}
\clearpage
\begin{exercise}
	Let $f : V \to W$ be differentiable. Let $\gamma : [0,1] \to V$ be differentiable. Then $f \circ \gamma$ is differentiable with
	\begin{align*}
		(f \circ \gamma)'(t)
		&=
		Df(\gamma'(t)) \cdot \gamma(t)
	\end{align*}
\end{exercise}
\begin{theorem}
	Let $U \subseteq V$ be open. Let $f : U \to W$ be differentiable with $Df = 0$. Then $f$ is locally constant.
\end{theorem}
\begin{corollary}
	Let $U \subseteq V$ be open and connected. Let $f : U \to W$ be differentiable with $Df = 0$. Then $f$ is constant.
\end{corollary}

\clearpage

\section{Higher Partial Derivatives}
\begin{importanttheorem}{Schwartz's Theorem}{}
	Let $f : U \subseteq V \to W$ be twice continuously partially differentiable. Then the order of partial derivation does not matter, i.e.
	\begin{align*}
		\partial_i \partial_j f = \partial_j \partial_i f
	\end{align*}
	for all $i,j \in [1..n]$.
\end{importanttheorem}
\begin{importanttheorem}{First-Order Taylor Expansion in multiple Variables}{}
	Let $f : U \subseteq V \to W$. Let $f$ be continuously differentiable. Let $x \in U$, $h \in V$ fulfill $[x,x+h] \subseteq V$. Then
	\begin{align*}
		f(x+h) = f(x) + Df(x)(h) + r(h),
	\end{align*}
	with $r \in o(\norm{h})$.
\end{importanttheorem}

\clearpage

\section{Differential Operators}
\chapter{The Two Queens of Multivariable Calculus}
\section{The Inverse Function Theorem}
\begin{importanttheorem}{Inverse Function Theorem}{}
	Let $U \subseteq \bR^n$ be open. Let $f : U \to \bR^n$ be continuously differentiable. Let $x \in U$ be a point such that $Df(x) : \bR^n \to \bR^n$ is invertible. Then there exists an open neighbourhood $V$ of $x$ such that:
	\begin{enumerate}
		\item $f(V)$ is an open neighbourhood of $f(x)$.
		\item $f|_V$ is a $C^1$-diffeomorphism, i.e. invertible with continuously differentiable inverse.
	\end{enumerate}
\end{importanttheorem}
In short: if a continuously differentiable function has an invertible differential at a given point, there exists an open neighbourhood around that point in which the function is invertible:
\begin{corollary}
	Let $U \subseteq \bR^n$ be open. Let $f : U \to \bR^n$ be continuously differentiable. Let $Df(x)$ be invertible for all $x \in \bR^n$. Then $f(U)$ is open.
\end{corollary}
\section{The Implicit Function Theorem}
\begin{importanttheorem}{Implicit Function Theorem}{}
	Let $U \subseteq \bR^m \times \bR^k$ be open. Let $f : U \to \bR^k$ be continuously differentiable.
	
	Let $(x,y) \in U$ with $f(x,y) = z$. Let $Df|_{\bR^k}(x,y) : \bR^k \to \bR^k$ be invertible. Then there exist
	\begin{enumerate}
		\item open neighbourhoods $V$ of $x$ and $W$ of $y$,
		\item and a continuously differentiable function $g : V \to W$,
	\end{enumerate}
	such that the set of values in $V \times W$ for which $f$ takes the value $z$ is the graph of $g$, i.e.
	\begin{align*}
		\lr{(v,w) \in V \times W \mid f(v,w) = z}
		=
		\lr{(v,g(v)) \mid v \in V}
	\end{align*}
\end{importanttheorem}
In short: under mild assumptions, the level sets of a continuously differentiable function are graphs of continuously differentiable functions:
\chapter{Multiple Integrals}
\begin{theorem}
	Let $K \subseteq \bR^n$ be compact. Let $f : K \times [a,b] \to \bR$ be continuous. Then the function
	\begin{align*}
		x \mapsto \int_a^b f(x,t) ~dt
	\end{align*}
	is continuous.
\end{theorem}
\begin{proof}
	$K \times [a,b]$ is a finite product of compact sets, i.e. compact.
\end{proof}
\chapter{Curves}
\chapter{Surfaces}